% GENERATED by scripts/gen_blueprint.py — do not edit.
\chapter{The Dictionary: every definition you need}
\label{ch:dictionary}

\begin{definition}[Digraphs and arcs ($u \longrightarrow v$)]
  \label{def:arc}
  \rocqok

\begin{verbatim}
HB.mixin Record HasArc V := { arc : rel V }.
Notation "u --> v" := (arc u v) (at level 30) : digraph_scope.
\end{verbatim}
\href{https://github.com/LLM4Rocq/digraph-theory/blob/main/theories/core/digraph.v\#L35}{Source: \code{theories/core/digraph.v:35}}

A digraph is a (finite, in this document) vertex type equipped with an
arc relation. The library roots its own hierarchy on a bare relation:



\code{rel V} is the type of Boolean binary relations on
\code{V}; \code{u --> v} is notation for
\code{arc u v}, ``there is an arc from $u$ to $v$''.
A \code{diGraphType} is a finite type bundled with such a
relation, with no further axioms.

\par\textbf{Why this is the right definition.} 
No axioms at this level means nothing is smuggled in: loops and
2-cycles are not yet excluded. They are excluded exactly where the
mathematics requires it --- by the oriented/tournament axioms below.
Isomorphism of digraphs is the evident notion (\ref{def:dgiso}).
\par

\end{definition}

\begin{definition}[Digraph isomorphism]
  \label{def:dgiso}
  \uses{def:arc}
  \rocqok

\begin{verbatim}
Definition dgiso (D1 D2 : diGraphType) : Prop :=
  exists f : D1 -> D2, bijective f /\ forall u v, (f u --> f v) = (u --> v).
\end{verbatim}
\href{https://github.com/LLM4Rocq/digraph-theory/blob/main/theories/core/digraph.v\#L97}{Source: \code{theories/core/digraph.v:97}}

A bijection between the vertex types that preserves and
reflects arcs ($f u \to f v$ iff $u \to v$, as Booleans).

\end{definition}

\begin{definition}[Oriented digraphs]
  \label{def:oriented}
  \uses{def:arc}
  \rocqok

\begin{verbatim}
HB.mixin Record DiGraph_IsOriented V of DiGraph V := {
  arc_irrefl : irreflexive (arc : rel V);
  arc_asymm  : forall u v : V, arc u v -> arc v u = false
}.
\end{verbatim}
\href{https://github.com/LLM4Rocq/digraph-theory/blob/main/theories/core/oriented.v\#L33}{Source: \code{theories/core/oriented.v:33}}

An \emph{oriented digraph} (an ``oriented graph'' in the combinatorics
literature) has no loops and no digons:



\code{arc\_irrefl} forbids loops; \code{arc\_asymm} says
$u \to v$ and $v \to u$ cannot both hold. The bundled structure is
\code{orientedDigraph}.

\end{definition}

\begin{definition}[Out-degree]
  \label{def:outdeg}
  \uses{def:arc}
  \rocqok

\begin{verbatim}
Definition outdeg v := #|[set w | v --> w]|.
\end{verbatim}
\href{https://github.com/LLM4Rocq/digraph-theory/blob/main/theories/core/oriented.v\#L47}{Source: \code{theories/core/oriented.v:47}}

The number of out-neighbours of $v$: $d^+(v) = |\{w : v \to
w\}|$. The path results state their hypothesis as
``\code{forall v, k <= outdeg v}'', i.e.\ minimum out-degree
$\delta^+ \ge k$.

\end{definition}

\begin{definition}[Tournaments]
  \label{def:tournament}
  \uses{def:oriented}
  \rocqok

\begin{verbatim}
HB.mixin Record Oriented_IsTournament V of Oriented V := {
  arc_total : forall u v : V, (u != v) = (arc u v) (+) (arc v u)
}.
\end{verbatim}
\href{https://github.com/LLM4Rocq/digraph-theory/blob/main/theories/core/tournament.v\#L32}{Source: \code{theories/core/tournament.v:32}}

A tournament orients every pair: for distinct $u,v$ exactly one of
$u \to v$, $v \to u$ holds.



\code{(+)} is exclusive or, so the axiom reads: $u \ne v$ iff
exactly one of the two arcs is present. Together with irreflexivity
(inherited from \code{Oriented}) this is precisely the textbook
definition. The bundled structure is \code{tournament}.

\par\textbf{Why this is the right definition.} 
On the diagonal the axiom specialises to $\mathrm{false} =
(u \to u) \oplus (u \to u)$, which holds trivially --- irreflexivity
is genuinely carried by the \code{Oriented} layer, and totality
by this one; no case is lost.
\par

\end{definition}

\begin{definition}[Transitive tournaments]
  \label{def:transb}
  \uses{def:arc}
  \rocqok

\begin{verbatim}
Definition transb := [forall u : T, [forall v : T, [forall w : T,
  (u --> v) ==> (v --> w) ==> (u --> w)]]].
\end{verbatim}
\href{https://github.com/LLM4Rocq/digraph-theory/blob/main/theories/core/tournament.v\#L112}{Source: \code{theories/core/tournament.v:112}}

The Boolean ``the arc relation is transitive''. Transitive
tournaments are exactly the linearly ordered ones; they are the
tournaments with $\bar\omega = 1$
(\code{omegabar\_transb}, \ref{res:omegabar_transb}).

\end{definition}

\begin{definition}[The directed triangle $C_3$]
  \label{def:C3}
  \uses{def:tournament}
  \rocqok

\begin{verbatim}
Definition C3 : Type := 'Z_3.

Lemma arcC3E (u v : C3) : (u --> v) = (v == u + 1 :> 'Z_3).
\end{verbatim}
\href{https://github.com/LLM4Rocq/digraph-theory/blob/main/theories/core/tournament.v\#L183}{Source: \code{theories/core/tournament.v:183}}

Vertices $\mathbb{Z}/3$, arcs $u \to u+1$: the 3-cycle
$0 \to 1 \to 2 \to 0$, the smallest non-transitive tournament.

\end{definition}

\begin{definition}[Vertex orders as permutations]
  \label{def:ltp}
  \rocqok

\begin{verbatim}
Definition ltp p u v := enum_rank (p u) < enum_rank (p v).

Definition realize : {perm V} := perm rfun_inj.

Lemma ltp_realizeE u v : ltp realize u v = r u v.
\end{verbatim}
\href{https://github.com/LLM4Rocq/digraph-theory/blob/main/theories/core/order.v\#L39}{Source: \code{theories/core/order.v:39}}

$\bar\omega$ quantifies over \emph{orderings of the vertex set}. The
library represents an ordering by a permutation $p$, compared through
the position of the enumeration:



\code{ltp p u v} reads ``$u$ precedes $v$ in the order named by
$p$''.

\par\textbf{Why this is the right definition.} 
Two directions need checking. (i) Every \code{ltp p} is an
irreflexive, transitive, total comparison --- a strict linear order on
the vertices. (ii) Conversely, \emph{every} strict linear order arises
this way: \code{realize} turns any irreflexive transitive total
relation \code{r} into a permutation, and
\code{ltp\_realizeE} states \code{ltp realize u v = r u v}.
So quantifying over \code{\{perm T\}} is quantifying over all
$|T|!$ vertex orders --- no order is missed, which is what makes the
\emph{minimum} in $\bar\omega$ honest.
\par

\end{definition}

\begin{definition}[The backedge graph]
  \label{def:backedge}
  \uses{def:tournament, def:ltp}
  \rocqok

\begin{verbatim}
Definition backedge_rel : rel T :=
  [rel u v | ltp p u v && (v --> u) || ltp p v u && (u --> v)].

Definition backedge : sgraph := SGraph backedge_sym backedge_irrefl.

Lemma backedgeE (u v : backedge) :
  (u -- v) = ltp p u v && (v --> u) || ltp p v u && (u --> v).
\end{verbatim}
\href{https://github.com/LLM4Rocq/digraph-theory/blob/main/theories/core/order.v\#L136}{Source: \code{theories/core/order.v:136}}

Fix a tournament $T$ and an order $p$. The backedge graph has an
(undirected) edge between $u$ and $v$ when the arc between them points
\emph{backwards} with respect to $p$:



\code{sgraph} is the simple-graph type of the
\texttt{coq-graph-theory} library, whose clique machinery
(\code{cliques}, \code{omega}) is reused unchanged;
\code{u -- v} is its edge relation.

\par\textbf{Why this is the right definition.} 
A set $K$ is a clique of the backedge graph iff, listing $K$ in
$p$-increasing order, every later vertex beats every earlier one ---
i.e.\ $K$ is a ``reverse-ordered'' subtournament under $p$. This is
the standard combinatorial reading of backedge cliques.
\par

\end{definition}

\begin{definition}[The tournament clique number $\bar\omega$]
  \label{def:omegabar}
  \uses{def:backedge}
  \rocqok

\begin{verbatim}
Definition omegab_at (T : tournament) (p : {perm T}) : nat :=
  ω([set: backedge p]).

Definition omegabar (T : tournament) : nat :=
  omegab_at [arg min_(p < (1%g : {perm T})) omegab_at p].
\end{verbatim}
\href{https://github.com/LLM4Rocq/digraph-theory/blob/main/theories/invariants/omegabar.v\#L29}{Source: \code{theories/invariants/omegabar.v:29}}

\code{omegab\_at p} is the clique number of the backedge graph of
the order $p$ (graph-theory's \code{omega} of the full vertex
set); \code{omegabar} takes the \emph{minimum} over all $p$
(\code{[arg min\_(...)]} is MathComp's minimiser over the finite
type \code{\{perm T\}}). The display notation in the source is
the omega-bar of the quoted statements.

\par\textbf{Why this is the right definition.} 
$\bar\omega(T) = \min_{\prec} \omega(\mathrm{BE}(T,\prec))$ over all
$|T|!$ linear orders $\prec$, by the realization argument of
\ref{def:ltp}. This matches Aboulker--Aubian--Charbit--Lopes
(arXiv:2310.04265, where the invariant measures the distance of $T$
from transitivity); $\bar\omega(T) = 1$ iff $T$ is transitive
(\ref{res:omegabar_transb}).
\par

\end{definition}

\begin{definition}[$k$-$\bar\omega$-critical tournaments]
  \label{def:kcritical}
  \uses{def:omegabar, def:subdel}
  \rocqok

\begin{verbatim}
Definition kcritical (k : nat) (T : tournament) : bool :=
  (ω̄(T) == k) && [forall v : T, ω̄(del_tournament v) == k.-1].
\end{verbatim}
\href{https://github.com/LLM4Rocq/digraph-theory/blob/main/theories/invariants/critical.v\#L21}{Source: \code{theories/invariants/critical.v:21}}

$T$ is $k$-critical when $\bar\omega(T) = k$ and deleting
\emph{any single} vertex drops $\bar\omega$ to $k-1$. The
\code{[forall v : T, ...]} ranges over all vertices;
\code{k.-1} is $k - 1$.

\end{definition}

\begin{definition}[Subtournaments and vertex deletion]
  \label{def:subdel}
  \uses{def:tournament}
  \rocqok

\begin{verbatim}
Definition sub_tournament (T : tournament) (S : {set T}) : tournament :=
  {x : T | x \in S} : tournament.

Definition del_tournament (T : tournament) (v : T) : tournament :=
  sub_tournament [set~ v].

Lemma sub_arcE (u v : {x | P x}) : (u --> v) = (val u --> val v).
\end{verbatim}
\href{https://github.com/LLM4Rocq/digraph-theory/blob/main/theories/core/tournament.v\#L171}{Source: \code{theories/core/tournament.v:171}}

For $S \subseteq V(T)$, \code{sub\_tournament S} is the
induced subtournament: its vertices are the elements of $S$ (a
dependent-pair type, whose first projection \code{val} is the
underlying vertex) and its arcs are inherited verbatim
(\code{sub\_arcE}). Deleting $v$ is the special case
$S = V \setminus \{v\}$ (\code{[set\textasciitilde{} v]}).

\end{definition}

\begin{definition}[Directed domination]
  \label{def:domination}
  \uses{def:tournament}
  \rocqok

\begin{verbatim}
Definition dominatesb (X : {set T}) : bool :=
  [forall v : T, (v \in X) || [exists x in X, x --> v]].

Definition domnum : nat := #|[arg min_(X < [set: T] | dominatesb X) #|X|]|.
\end{verbatim}
\href{https://github.com/LLM4Rocq/digraph-theory/blob/main/theories/invariants/domination.v\#L28}{Source: \code{theories/invariants/domination.v:28}}

$u$ dominates $v$ if $u = v$ or $u \to v$;
$\mathrm{dom}(T)$ is the least size of a set dominating every vertex.
Used in the general bound $\mathrm{dom}(T) \le \bar\omega(T)$
(\ref{res:domnum_le_omegabar}).

\end{definition}

\begin{definition}[Automorphisms]
  \label{def:aut}
  \uses{def:arc}
  \rocqok

\begin{verbatim}
Definition autb p : bool :=
  [forall u : D, [forall v : D, arc (p u) (p v) == arc u v]].

Definition dgaut : {set {perm D}} := [set p | autb p].
\end{verbatim}
\href{https://github.com/LLM4Rocq/digraph-theory/blob/main/theories/constructions/automorphism.v\#L25}{Source: \code{theories/constructions/automorphism.v:25}}

A permutation of the vertices preserving arcs in both
directions; \code{dgaut} is the set of all of them (a group).

\end{definition}

\begin{definition}[Vertex-transitivity]
  \label{def:vt}
  \uses{def:aut}
  \rocqok

\begin{verbatim}
Definition vertex_transitiveb : bool :=
  [forall u : D, [forall v : D,
    [exists p : {perm D}, (p \in dgaut) && (p u == v)]]].
\end{verbatim}
\href{https://github.com/LLM4Rocq/digraph-theory/blob/main/theories/constructions/automorphism.v\#L55}{Source: \code{theories/constructions/automorphism.v:55}}

For every pair of vertices some automorphism carries one to
the other. All three critical families here are vertex-transitive ---
which is why checking one deletion suffices
(\ref{res:vt_kcritical}).

\end{definition}

\begin{definition}[Lexicographic substitution $S[H]$]
  \label{def:lexprod}
  \uses{def:tournament}
  \rocqok

\begin{verbatim}
Definition lexprod : Type := (D1 * D2)%type.

Lemma lexprod_arcE (u v : lexprod) :
  (u --> v) = (u.1 --> v.1) || (u.1 == v.1) && (u.2 --> v.2).

Definition lexprod_tournament (T1 T2 : tournament) : tournament :=
  lexprod T1 T2 : tournament.
\end{verbatim}
\href{https://github.com/LLM4Rocq/digraph-theory/blob/main/theories/constructions/product.v\#L25}{Source: \code{theories/constructions/product.v:25}}

Vertices are pairs $(s,h)$; between different $S$-blocks the
arc follows the $S$-arc of the first coordinates, inside a block it
follows $H$ (\code{lexprod\_arcE} is the characterising equation:
the arc is \code{(u.1 --> v.1) || (u.1 == v.1) \&\& (u.2 --> v.2)},
and for tournaments the first disjunct can only fire when
$u.1 \neq v.1$). This is the standard substitution/composition
$S[H]$: every block is a copy of $H$, blocks are arranged as $S$.

\end{definition}

\begin{definition}[Cayley digraphs]
  \label{def:cayley}
  \uses{def:arc}
  \rocqok

\begin{verbatim}
Definition cayley (gT : finGroupType) (A : {set gT}) : Type := gT.

Lemma cayley_arcE (x y : cayley A) : (x --> y) = ((x^-1 * y)%g \in A).
\end{verbatim}
\href{https://github.com/LLM4Rocq/digraph-theory/blob/main/theories/constructions/cayley.v\#L30}{Source: \code{theories/constructions/cayley.v:30}}

For a finite group $G$ and connection set $A \subseteq G$:
vertices $G$, arc $x \to y$ iff $x^{-1}y \in A$. Over
$\mathbb{Z}/n$ (written additively) this is the circulant: $x \to y$
iff $y - x \in A$.

\end{definition}

\begin{definition}[The circulant tournament $\mathrm{AC}_n$]
  \label{def:AC}
  \uses{def:cayley, def:tournament}
  \rocqok

\begin{verbatim}
Definition ACset : {set 'Z_n} :=
  [set z : 'Z_n | ((0 < val z < m) || (val z == m.+1))%N].

Definition AC : tournament := cayley ACset : tournament.

Lemma AC_arcE (x y : AC) : (x --> y) = (y - x \in ACset).
\end{verbatim}
\href{https://github.com/LLM4Rocq/digraph-theory/blob/main/theories/constructions/circulant.v\#L63}{Source: \code{theories/constructions/circulant.v:63}}

The paper's witness platform: for odd $n = 2m+1$, the circulant with
connection set $g = \{1, \dots, m-1\} \cup \{m+1\}$:



(\code{val z} is the representative of $z \in \mathbb{Z}/n$ in
$\{0,\dots,n{-}1\}$; the set comprehension says $0 < \mathrm{val}\, z
< m$ or $\mathrm{val}\, z = m+1$.) Since $g$ contains exactly one of
$\{z, -z\}$ for each $z \ne 0$, $\mathrm{AC}_n$ is a tournament; as a
Cayley digraph it is vertex-transitive. Throughout, $m$ is
\code{m'.+1} and $n$ is \code{m.*2.+1}
(\ref{ch:primer}). In statements quoted from section contexts, the
local abbreviations \code{ACm}, \code{m}, \code{n}
denote \code{AC m'}, \code{m'.+1}, \code{m.*2.+1}.

\end{definition}

\begin{definition}[The $k=4$ family $\mathrm{AC}_n[C_3]$]
  \label{def:T4family}
  \uses{def:lexprod, def:AC, def:C3}
  \rocqok


\code{T4} abbreviates
\code{lexprod\_tournament ACm (C3 : tournament)}: the
substitution of the directed triangle into $\mathrm{AC}_n$
(\ref{def:lexprod}, \ref{def:AC}, \ref{def:C3}), a vertex-transitive tournament on
$3n$ vertices.

\end{definition}

\begin{definition}[The $k=5$ family $\mathrm{AC}_n[\mathrm{AC}_n]$]
  \label{def:T5family}
  \uses{def:lexprod, def:AC}
  \rocqok


\code{T5} abbreviates \code{lexprod\_tournament ACm ACm}:
$\mathrm{AC}_n$ substituted into itself, a vertex-transitive
tournament on $n^2$ vertices.

\end{definition}

\begin{definition}[Directed simple paths]
  \label{def:dipath}
  \uses{def:arc}
  \rocqok

\begin{verbatim}
Definition dipath x s := path arc x s && uniq (x :: s).
\end{verbatim}
\href{https://github.com/LLM4Rocq/digraph-theory/blob/main/theories/core/dipath.v\#L36}{Source: \code{theories/core/dipath.v:36}}

\code{path arc x s} (MathComp) says consecutive
elements of \code{x :: s} are joined by arcs;
\code{uniq} makes the path simple. Its \emph{length} (number of
arcs) is \code{size s}.

\end{definition}

\begin{definition}[The longest-path length $\ell(D)$]
  \label{def:ell}
  \uses{def:dipath}
  \rocqok

\begin{verbatim}
Definition has_dipath k := [exists x : D, exists t : k.-tuple D, dipath x t].

Definition ell := \max_(k < #|D| | has_dipath k) (k : nat).
\end{verbatim}
\href{https://github.com/LLM4Rocq/digraph-theory/blob/main/theories/core/dipath.v\#L144}{Source: \code{theories/core/dipath.v:144}}

$\ell(D)$ is the largest $k$ for which a directed simple
path with $k$ arcs exists (\code{\max\_(...)} over the finite
range; a simple path has at most $|D| - 1 < |D|$ arcs, so the bounded
maximum is the true maximum). The display notation in the source is the script ell of
the quoted statements.

\end{definition}

\begin{definition}[Strong connectivity]
  \label{def:strong}
  \uses{def:arc}
  \rocqok

\begin{verbatim}
Definition strongb := [forall x : D, forall y : D, connect arc x y].
\end{verbatim}
\href{https://github.com/LLM4Rocq/digraph-theory/blob/main/theories/invariants/strong.v\#L37}{Source: \code{theories/invariants/strong.v:37}}

Every vertex reaches every other by a directed walk
(\code{connect} is MathComp's reflexive-transitive closure).

\end{definition}

